\chapter{一台能运行程序的机器为什么仍需要操作系统}

\begin{statebox}
\textbf{逻辑起点。} 机器只有复位、时钟、处理器、易失内存、非易失启动存储、最小互连和
一项待运行任务。当前还没有 kernel、task、scheduler、system call 或文件。第一项任务
要求把四个整数相加，并把结果 \(21\) 写到约定地址；每个后续结构都必须由这台机器已经
暴露的具体失败推导出来。
\end{statebox}

\inputsamplechapter{chapters18/evolution/ch01-one-task}
\inputsamplechapter{chapters18/evolution/ch02-automatic-loading}
\inputsamplechapter{chapters18/evolution/ch02-resident-coordinator}
\inputsamplechapter{chapters18/evolution/ch03-forced-return}

\begin{problemframe}
\textbf{本章得到的机器。} 固件能够装入任务；常驻协调代码能够保存和恢复现场；计时器事件
使协调代码能从不合作的任务手中重新获得处理器。尚未解决的是：如何把每项执行流的长期状态
组织成可扩展对象，如何表示可运行与等待关系，以及多核条件下由谁拥有这些队列。下一章只
处理这组执行状态，不提前展开用户请求、地址空间或文件。
\end{problemframe}
